\documentclass[12pt,a4paper]{article}

\usepackage{cmap}
\usepackage[T2A]{fontenc}
\usepackage[utf8]{inputenc}
\usepackage[russian]{babel}
\usepackage{amsmath,amssymb}
\usepackage[a4paper,margin=2cm]{geometry}

\setlength{\parindent}{0pt}
\setlength{\parskip}{0.45em}
\sloppy

\begin{document}

\begin{center}
{\Large\bfseries Билет 11: вопросы для проверки знаний}
\end{center}

\noindent\fbox{%
  \begin{minipage}{\dimexpr\linewidth-2\fboxsep-2\fboxrule\relax}
  \normalfont
Равномерная сходимость функциональных последовательностей и рядов.
Критерий Коши равномерной сходимости. Непрерывность суммы равномерно
сходящегося ряда из непрерывных функций. Интегрирование и дифференцирование
функциональных последовательностей и рядов. Признаки Вейерштрасса, Дирихле и
Абеля равномерной сходимости функциональных рядов.
  \end{minipage}%
}

\section*{Вопросы на понимание}

\begin{enumerate}
\item Чем равномерная сходимость \(f_n \to f\) на множестве \(X\) сильнее поточечной сходимости?
\item Что означает фраза <<номер \(N\) выбирается один для всех \(x\)>> в определении равномерной сходимости?
\item Почему последовательность \(x^n\) ведет себя по-разному на \([0,1]\) и на \([0,a]\) при \(0<a<1\)?
\item Почему для функционального ряда равномерную сходимость проверяют по частичным суммам, а не только по отдельным членам \(u_n(x)\)?
\item В критерии Коши для функционального ряда зачем оценивать все конечные хвосты ряда, а не только один следующий член?
\item Почему условие \(u_n(x)\rightrightarrows 0\) необходимо, но само по себе не гарантирует равномерную сходимость ряда?
\item Что может пойти не так с непрерывностью суммы ряда из непрерывных функций, если сходимость только поточечная?
\item Почему при почленном интегрировании на \([a,b]\) важна равномерная оценка остатка, а не просто сходимость в каждой точке?
\item Почему равномерная сходимость самого ряда обычно не дает права дифференцировать его почленно?
\item В условии о дифференцировании зачем нужна сходимость значений \(f_n(x_0)\) или соответствующего числового ряда хотя бы в одной точке?
\item Почему в признаке Вейерштрасса числовая мажоранта \(a_k\) должна не зависеть от \(x\)?
\item Какую разную роль играют условия <<равномерно ограничены частичные суммы>> и <<\(b_k(x)\rightrightarrows 0\)>> в признаке Дирихле?
\end{enumerate}

\section*{Источники для сверки}

\texttt{target/answer\_question\_11.tex};
\texttt{IvGE\_1.pdf}, стр. 281--299, 303.

\end{document}
